\section{Notation, prerequisites and quoted results}\label{sec:prelim}%
 

\subsection{General notation}
Throughout this manuscript, unless otherwise specified, we adopt the following notation and conventions.
\par
All groups considered are finite, and $G$ and $H$ always denote finite groups. We let $k$ be an algebraically closed field of characteristic~$p>0$.  From Section~\ref{sec:blocksWithDihedralDefectGroups} on, we assume that $p=2$.
When character-theoretic  arguments are needed, we let  $\cO$ be a complete discrete valuation ring with residue field $\cO/J(\cO)=k$ and a quotient field $K$ of characteristic zero, which is
assumed to be large enough for the finite groups under consideration and their subgroups. 
\par 
For $R\in\{k,\cO\}$ all modules over a given $R$-algebra $\bA$ that is free of finite rank over \(R\), are assumed to be finitely generated left
%\todo[color=yellow]{(!)We use left modules.} 
modules and we denote by $\bA-\mathsf{mod}$ the corresponding module category.  
The set of all blocks of $\bA$ (i.e. its indecomposable algebra factors) is denoted $\Bl(\bA)$.
We write $\bB_0(RG)$ for the principal block of $RG$. 
Given $\bB\in\Bl(RG)$, we write  $e_{\bB}$  for the block idempotent of $RG$ corresponding to $\bB$, and 
if $D\leq G$ is a defect group of $\bB$, then the integer $d$ such that $|D|=p^{d}$ is called the defect of $\bB$. We say that $\bB$ is a block of \emph{maximal defect} provided its defect groups are the Sylow $p$-subgroups of $G$, i.e.  $|G|_p=p^d$, and 
%$\Bl(kG\mid D):=\{\bB\in\Bl(kG)\mid D\text{ is a defect group of }\bB\}$, 
we let $\Bl^{\mathrm{max}}(RG)$ denote  the subset of $\Bl(RG)$ consisting of all blocks of maximal defect of $RG$. If $N\unlhd G$ and $\bc\in\Bl(RN)$, then  $\Bl(RG\mid \bc):=\{\bB\in\Bl(RG)\mid \bB\text{ covers }\bc\}$.


\par
The canonical algebra homomorphism $\cO G\longrightarrow kG$ induces a defect group preserving bijection between $\Bl(\cO G)$ and $\Bl(kG)$, mapping a block $\bB\in \Bl(\cO G)$ to $k\otimes_{\cO}\bB$. In this context, we write $\Irr(\bB)$,  and also simply $\Irr(k\otimes_{\cO}\bB)$ when working over $k$, 
%$\IBr(\bB)$, respectively,
for the sets of all
irreducible 
%and Brauer 
$K$-characters 
of $G$ belonging to $\bB$. We write $\Irr(G)$ for the set of all irreducible $K$-characters of~$G$.
\par
In contrast, if $\bB\in\Bl(kG)$, then $\Irr_k(\bB)$ denotes the set of all simple $\bB$-modules up to isomorphism. 
%If $S$ is a simple $\bB$-module affording the Brauer character~$\varphi$, to simplify, we may sometimes identify $S$ and $\varphi$ with $d:=\dim_k(S)$ and, by abuse of notation, we write $S=d$ or $\varphi=d$.
We denote by $k_G$, or simply $k$,  the trivial $kG$-module. 
For $H\leq G$ we denote by $\Res^G_H$ (resp.~$\Ind_H^G$) the restriction (resp. the induction) of $kG$-modules and characters. Similarly, for $N\unlhd G$, we denote by $\Inf_{G/N}^{G}$ the inflation of $kG$-modules and characters. Given a $p$-subgroup $Q\leq G$ and a $kG$-module $M$  we denote by $M[Q]$ the Brauer construction of $M$ with respect to $Q$, and by $\Br^G_Q$ (or simply $\Br_Q$ when clear from the context)  the Brauer homomorphism  with respect to $Q$, i.e. the surjective algebra homomorphism $\Br^G_Q: (kG)^Q\longrightarrow kC_G(Q)$
mapping $C_G(Q)$ to itself by the identity map and mapping all the other group
elements to zero. The projective cover of a simple $kG$-module $S$ is written $P(S)$. If $S_1,S_2$ are simple $kG$-modules, then $U(S_1,S_2)$ denotes a uniserial $kG$-module with head $S_1$ and socle $S_2$. 
\par
Again, given $R\in\{\mathcal{O},k\}$, an $RG$-lattice $M$ is a \emph{permutation} $RG$-lattice if it admits a $G$-invariant $R$-basis, whereas it is a $p$-\emph{permutation} $RG$-lattice if it is a direct summand of a permutation $RG$-lattice. Indecomposable $p$-permutation $RG$-lattices are also called \emph{trivial source} $RG$-lattices. 
It is well known that any $p$-permutation $kG$-module $M$ is liftable to an $\cO G$-lattice. In general, such modules afford several lifts, but there is a unique one amongst these which is a $p$-permutation $\cO G$-lattice. We denote this lift by $\wh{M}$ and we let $\chi_{\wh{M}}$ be the character afforded by $K\otimes_{\cO}\wh{M}$.
For any $p$-subgroup $Q$ of $G$, the Scott module with vertex $Q$, written $\Sc(G,Q)$, is always a trivial source module. 
\par
Given an integer $n\geq 2$, 
$C_n$ denotes the cyclic group of order $n$, $\fA_n$ and $\fS_n$ the alternating and the symmetric group on $n$ letters, respectively, and $D_{2^n}$ the dihedral group of order $2^n$. 
Furthermore,   ${\mathrm{Syl}}_p(G)$ denotes  the set of all Sylow $p$-subgroups of $G$, and for $P\in\Syl_p(G)$ we let $\mathcal F_P(G)$ be the fusion system of $G$ on $P$. 
If $\varphi:Q\longrightarrow P$ is an isomorphism between subgroups $Q\leq H$ and $P\leq G$, then $\Delta(P,\varphi,Q):=\{(\varphi(y),y)\in P\times Q\mid y\in Q\}$
is a twisted diagonal subgroup of $G\times H$, and if $G=H$ and $P=Q$ then $\Delta(P):=\Delta(P,\Id,P)$ is the diagonal embedding of $P$ in $G\times G$. Further, we denote the (external) central product of two finite groups $G_1$ and $G_2$ identified along a group isomorphism $\varphi:Z_1\longrightarrow Z_2$ with $Z_i\leq Z(G_i)$ for $i=1,2$ by $G_1\ast G_2$, so 
$G_1\ast G_2=(G_1\times G_2)/Z$ with $Z=\{(z,\varphi(z)^{-1})\mid z\in Z_1\}\leq Z(G_1\times G_2)$.
\par
For further standard notation and terminology, we refer the reader to \cite{LinckelmannVol1,LinckelmannVol2,NagaoTsushimaBook,ThevenazBook}.

\vspace{2mm}
\subsection{Definitions: Morita equivalences, stable equivalences of Morita type, splendid Morita equivalences and endo-permutation source Morita equivalences}\label{ssec:Splendid_Morita_equivalences}{\leavevmode}%

\noindent Let $\bA$ and $\bB$ be finite-dimensional $R$-algebras with $R\in\{\cO,k\}$. We recall that a pair $(M,N)$ consisting of an $(\bA,\bB)$-bimodule $M$ and of a $(\bB,\bA)$-bimodule $N$ is said to \emph{induce a Morita
equivalence between $\bA$ and $\bB$} if $M$ and $N$ are finitely generated projective as left
and right modules, and if there are isomorphisms $M\otimes_{\bB} N\cong \bA$ and $N\otimes_{\bA}M\cong \bB$ of $(\bA,\bA)$-bimodules and $(\bB,\bB)$-bimodules respectively. In this case, the functors $M\otimes_{\bB}–:\bB-\mathrm{mod}\longrightarrow \bA-\mathsf{mod}$ and $N\otimes_{\bA}–:\bA-\mathrm{mod}\longrightarrow \bB-\mathsf{mod}$ are inverse equivalences of categories, and $\bA$ and $\bB$ are called \emph{Morita equivalent}.  Moreover, as bimodules, $N\cong  \Hom_{\bA}(M,\bA)$ and $M\cong \Hom_{\bB}(N,\bB)$, so we often simply say that $M$ \emph{induces a Morita
equivalence between $\bA$ and $\bB$}.  See~\cite[Theorem~2.8.2]{LinckelmannVol1}. 
Under the weaker assumptions that $M\otimes_{\bB} N\cong \bA\oplus(\text{projective $(\bA,\bA)$-bimodule})$ and $N\otimes_{\bA}M\cong \bB\oplus(\text{projective $(\bB,\bB)$-bimodule})$,  it is said that the pair $(M,N)$ induces a \emph{stable equivalence of Morita type} between $\bA$ and $\bB$. 

\par
Assume now that $G$ and $H$ are finite groups, $\bA\in\Bl(RG)$, and $\bB\in\Bl(RH)$. Then, the $(\bA,\bB)$-bimodule $M$ can be regarded as a left $R[G\times H]$-module (via the left action given by $(g,h)\cdot m:=g\cdot m\cdot h^{-1}$ for any $(g,h)\in G\times H$ and any $m\in M$) and the $(\bB,\bA)$-bimodule $N$ as a left $R[H\times G]$-module. The Morita equivalence is called a \emph{splendid Morita equivalence} (respectively, an \emph{endo-permutation source Morita equivalence}) if $M$ and $N$, seen as left modules, are trivial source modules (respectively, have an endo-permutation source). Clearly, any splendid Morita equivalence is an endo-permutation source  Morita equivalence. If there is an endo-permutation source Morita equivalence between $\bA$ and $\bB$, then they  have isomorphic defect groups. After identifying defect groups of $\bA$ and $\bB$ with a $p$-group $D$, the modules $M$ and $N$ have vertex $\Delta D$. A splendid Morita equivalence is sometimes also called a \emph{Puig equivalence} or a \emph{source-algebra equivalence}, since it is equivalent to the existence of an isomorphism between source algebras of $\bA$ and $\bB$ as interior $D$-algebras. See \cite[Proposition~9.7.1, Remark~9.7.2, Definition~9.7.3,  Theorem~9.7.4 and Theorem 9.11.2]{LinckelmannVol2}.
\vspace{2mm}
\subsection{Blocks of direct and central products}\label{ssec:blocks_of_central_producrs}
Let $G$ and $H$ be finite groups as in the previous subsection. 
Recall that there is a canonical $R$-algebra isomorphism $R[G\times H]\longrightarrow RG\otimes_R RH$, mapping $(g,h)\in G\times H$ to $g\otimes h$. Identifying $R[G\times H]$ and $RG\otimes_R RH$ via this isomorphism, each block $\bB \in\Bl(R[G\times H])$ can be written as $\bB=\bB_G\otimes_R\bB_H\in \Bl(RG\otimes_R RH)$ with $\bB_G\in\Bl(RG)$ and $\bB_H\in\Bl(RH)$. 
\par
More generally, if $G\ast H$ is a central product of $G$ and $H$ identified along a central subgroup of $p'$-order, then any block $\bB \in\Bl(R[G\ast H])$ is isomorphic to $\bB_G\otimes_R\bB_H$ where $\bB_G\in\Bl(RG)$ is the unique block of $RG$ covered by $\bB$ and $\bB_H\in\Bl(RH)$ is the unique block of $RH$ covered by $\bB$. See \cite[Proposition~1.5]{Duvel04} (or also \cite[Lemma~7.5]{Sambale14}). More precisely, the block $\bB_G\otimes_R\bB_H$, viewed as a block of \(R[G\times H]\), maps onto $\bB$ under the canonical map $R[G\times H]\to R[G\ast H]$. By
\cite[Lemma~1.4]{Duvel04}, this map induces an algebra isomorphism
\[
        \bB_G\otimes_R\bB_H\cong \bB.
\]
Since the kernel of the canonical homomorphism $G\times H\to G\ast H$ is a $p'$-group, the quotient map identifies the corresponding defect groups. If $D$ denotes the common defect group, then the above map, being induced by a group homomorphism, is compatible with the corresponding interior $D$-algebra structures. Thus, $\bB_G\otimes_R\bB_H$ and $\bB$ have isomorphic source algebras as interior $D$-algebras. In particular, they are splendidly Morita equivalent. 



\begin{lem}\label{lem:SplendidMoritaEquivalence_direct_and_central_product}%
Let $G\ast H$ be a central product of finite groups identified along a central subgroup of $p'$-order (possibly trivial).
Let $\bB \in\Bl(R[G\ast H])$ and let $\bB_G\in\Bl(RG)$ and $\bB_H\in\Bl(RH)$ be the unique blocks covered by $\bB$, as above. If $\bB_G$ is a block of defect zero, then $\bB$ and $\bB_H$ are splendidly Morita equivalent. 
\end{lem}

\begin{proof}
By the above discussions, the blocks $\bB$ and $\bB_G\otimes_R\bB_H$ are splendidly Morita equivalent. By \cite[p.\,341, line -9]{ThevenazBook}, $R$ is a source algebra of $\bB_G$, thus the assertion follows from \cite[Lemma 2.6(i)]{KMN11}.
\end{proof}






\vspace{2mm}
\subsection{Groups with dihedral Sylow \texorpdfstring{$2$}{2}-subgroups}\label{ssec:dihedralSylow2}%
Let $G$ be a finite group with a dihedral Sylow $2$-subgroup $P$ of order $2^n\geq 8$. Then, Gorenstein and Walter~\cite[Theorem~1]{GorensteinWalter} proved 
that $\overline{G}:=G/O_{2'}(G)$ is isomorphic to $P$, the alternating group $\fA_7$,  or a subgroup of $\PGammaL_2(q)$ containing $\PSL_2(q)$, 
where $q:=r^f$ is a power of an odd prime $r$ and $f\geq 1$ is an integer. Then,  
$\PGammaL_2(q)\cong \PGL_2(q)\rtimes 
{\mathrm{Gal}}(\mathbb F_{q}/\mathbb F_{r})$,  
where ${\mathrm{Gal}}(\mathbb F_{q}/\mathbb F_{r})$  
is cyclic of order $f$. Thus, for the purposes of our computations, we split the above list of groups into further subcases as follows:
\begin{enumerate}[label={\bf(D\arabic*)}, leftmargin=12mm]
   \item $\overline{G}\cong D_{2^n}$;
   \item $\overline{G}\cong \fA_7$;
   %\item if $\overline{G}\cong {\mathrm{PSL}}(2,9) \cong \mathfrak A_6$; 
   %\item if $\overline{G}\cong \PGL(2,9)\cong \mathfrak A_6.2_2$;
   %\item if $\overline{G}\cong \PSL(2,q)$,  where  $q$ is a power of an odd prime with $q\neq 9$, and $q\equiv\pm 1\pmod{8}$;
   %\item if $\overline{G}\cong \PGL(2,q)$,  where $q$ is a power of an odd prime with $q\neq 9$;
   \item  $\overline{G}\cong \PSL_2(q)\rtimes C_d$,  where $q$ is a power of an odd prime such that $(q-1)_2=2^{n}$ or $(q+1)_2=2^{n}$, $d$ is odd and $d\,|\,f$;
   \item $\overline{G}\cong \PGL_2(q)\rtimes C_d$,  where $q$ is a power of an odd prime such that $(q-1)_2=2^{n-1}$ or $(q+1)_2=2^{n-1}$, $d$ is odd and $d\,|\,f$.
\end{enumerate}
Notice that Case {\bf(D3)} happens provided $q\equiv\pm 1\pmod{8}$, whereas Case {\bf(D4)} happens when  $q\equiv\pm 1\pmod{4}$. In both cases, $d$ is odd. %, e.g. by\cite[Chapter~6~(8.9)]{SuzukiGroupTheoryII}. 
Moreover, $\PSL_2(9)\cong \fA_6$ and $\PGL_2(9)\cong \fA_6.2_2$ in the ATLAS notation. \\

\noindent Next, recall that the Schur multiplier of a finite group $G$ is the abelian group ${\mathrm{H}}^2(G,\mathbb{C}^\times)$ and there is an isomorphism 
\[
    {\mathrm{H}}^2(G,k^\times)\cong {\mathrm{H}}^2(G,\mathbb{C}^\times)_{2^\prime} \tag{$\ast$}
\]
(see e.g. \cite[2.1.14 Proposition]{KarpilovskySchurMultBook}). For the groups above, this group is as follows. 


\begin{lem}\label{lem:SchurMult_2prime}
Assume $p=2$. Let $G$ be a finite group with a dihedral Sylow $2$-subgroup of order $2^n$ with $n\geq 3$ and assume that $O_{2^\prime}(G)=1$. Then, we have 
\[
{\mathrm{H}}^2(G, k^\times)\cong  \begin{cases}
    C_3  & \text{ if } \,  G
\in \{\mathfrak{A}_7, \PSL_2(9)\}, \\
    1  & \text{otherwise}.
\end{cases}
\]
Moreover, if $G
\in \{\mathfrak{A}_7, \PSL_2(9)\}$ and $Z$ is a finite cyclic group of odd order on which $G$ acts trivially and $3\mid |Z|$, then 
\[
{\mathrm{H}}^2(G, Z)\cong {\mathrm{H}}^2(G, k^\times)\cong C_3\,.
\]
\end{lem}

\begin{proof} 
We go through the cases {\bf (D1)} to {\bf (D4)}.  First, if $G=D_{2^n}$, then ${\mathrm{H}}^2(G, k^\times)\cong 1$ as $G$ is a $2$-group; see e.g. \cite[Chapter 2, Theorem~7.4]{NagaoTsushimaBook}.
\par
If $G\in\{\fA_7,\PSL_2(9)\}$, then we read from the ATLAS \cite[p.10 and p.4 resp.]{ATLAS} that ${\mathrm{H}}^2(G,\mathbb{C}^\times)\cong C_6$, so ${\mathrm{H}}^2(G, k^\times)\cong C_3$ by ($\ast$). 
\par
If {\bf (D3)} holds with $d=1$, $G=\PSL_2(q)$  and $q\neq 9$, then ${\mathrm{H}}^2(G,\mathbb{C}^\times)\cong C_2$  by \cite[Chapter~3, \S3 with Table~5]{ATLAS}, hence  ${\mathrm{H}}^2(G, k^\times)\cong 1$ by ($\ast$).
\par
Next, if $G\cong \PGL_2(3)\cong\mathfrak S_4$, then we have ${\mathrm{H}}^2(G, \mathbb{C}^\times)\cong C_2$
 by \cite[2.12.3 Theorem]{KarpilovskySchurMultBook}, hence ${\mathrm{H}}^2(G, k^\times)\cong 1$ by ($\ast$). If $G=\PGL_2(9)$, then we read from the ATLAS \cite[bottom of p.5 together with page xxii]{ATLAS} and ($\ast$) that ${\mathrm{H}}^2(G, k^\times)\cong 1$. 
\par
We go through the remaining groups using \cite[Theorem 3.1(c)]{Jones74}, which states that if $N\unlhd G$ is such that $G/N$ is cyclic, then 
$|{\mathrm{H}}^2(G,\mathbb{C}^\times)|$ divides $|{\mathrm{H}}^2(N,\mathbb{C}^\times)|\cdot |N/[N,N]|$. In all cases, we take $N:=\PSL_2(q)$, which is perfect. If $G$ is of type {\bf (D3)} with $d\neq 1$, then $G/N\cong C_d$, and certainly $q>9$.  If $G$ is of type {\bf (D4)} and $q\notin\{3,9\}$, then $G/N\cong C_{2d}$. In both cases, we obtain that 
\[
|{\mathrm{H}}^2(G,\mathbb{C}^\times)|\,\big|\,|{\mathrm{H}}^2(N,\mathbb{C}^\times)|\cdot |N/[N,N]|=2\cdot 1=2\,,
\]
thus ${\mathrm{H}}^2(G, k^\times)\cong 1$ by ($\ast$). This proves the first claim. The second claim follows from the Universal Coefficient Theorem, which reads  
\[
{\mathrm{H}}^2(G, Z) \cong \Hom_{\IZ}({\mathrm{H}}_2(G, \mathbb{Z}), Z) \oplus \Ext^1_{\mathbb{Z}}({\mathrm{H}}_1(G, \mathbb{Z}), Z)\,.
\]
(See e.g. \cite[Exercise~6.1.5]{WeibelBook}.)
Since $\fA_7$ and $\PSL_2(9)$ are perfect, the Ext part vanishes, because  ${\mathrm{H}}_1(G, \mathbb{Z})\cong G/[G,G]$. Moreover, $ {\mathrm{H}}_2(G,\mathbb{Z})\cong {\mathrm{H}}^2(G,\mathbb{C}^\times)$ and we obtain ${\mathrm{H}}^2(G, Z)\cong \Hom_{\IZ}(C_3, Z)\cong C_3$. 
\end{proof}



\begin{comment}
{
\color{blue}
\% COMMENTS
\begin{rem}
 The (equivalence classes of) central extensions $1\longrightarrow Z\longrightarrow \widehat{G}\longrightarrow G\longrightarrow 1$ of a finite group $G$ by a finite abelian group $Z$ are parametrized by ${\mathrm{H}}^2(G, Z)$, where $Z$ is seen as a trivial $G$-module. They are linked to the Schur multiplier through the Universal Coefficient Theorem, namely  
\[
{\mathrm{H}}^2(G, Z) \cong \Hom_{\IZ}({\mathrm{H}}_2(G, \mathbb{Z}), Z) \oplus \text{Ext}^1_{\mathbb{Z}}({\mathrm{H}}_1(G, \mathbb{Z}), Z)\,,
\]
where ${\mathrm{H}}^2(G,\mathbb{C}^\times)\cong {\mathrm{H}}_2(G,\mathbb{Z})$ and  ${\mathrm{H}}_1(G, \mathbb{Z})\cong G/[G,G]$ is the abelianization of $G$. (See e.g. \cite[Exercise 3 of  III.\S3]{BROWN}). 
\par
Assume now that $Z$ is a cyclic group of odd order and $G/[G,G]$ is a $2$-group, then 
${\mathrm{H}}^2(G, Z) \cong \Hom_{\mathbb{Z}}({\mathrm{H}}^2(G, \mathbb{C}^{\times}), Z)$, which is trivial provided ${\mathrm{H}}^2(G, k^{\times})$ is trivial and so $\widehat{G}\cong Z\times G$ is the trivial central extension.
\end{rem}
}
\end{comment}


\vspace{2mm}
\subsection{Endo-permutation modules and the Dade group}\label{subsec:endo-perm}%
Let $P$ denote a $p$-group. A $kP$-module $M$ is called \emph{endo-permutation} if $\End_k(M)$ is a permutation $kP$-module. Moreover, such an endo-permutation $kP$-module $M$ is called \emph{capped} provided it possesses an indecomposable direct summand with vertex $P$, denoted $\mathrm{cap}(M)$ and called its \emph{cap}. In this case, $\mathrm{cap}(M)$ is unique up to isomorphism, although it may appear with multiplicity greater than one, so $M=\mathrm{cap}(M)^a\oplus X$ with $a\in\mathbb{Z}_{>0}$ and $X$ is a $kP$-submodule of $M$, the indecomposable direct summands of which have vertices of order strictly smaller than $|P|$.   
This allows us to define an equivalence relation: two capped endo-permutation $kP$-modules are equivalent if and only if their caps are isomorphic. The \emph{Dade group} of $P$, denoted $\bD_k(P)$, is the resulting set of equivalence classes (of capped endo-permutation $kP$-modules) with group law given by the tensor product $\otimes_k$\,. By a result of Puig, $\bD_k(P)$ is a finitely generated abelian group. We denote its torsion part by $\bD_k^{tors}(P)$. For a detailed introduction to endo-permutation modules and their properties, we refer the reader to \cite{ThevenazSurvey}. \\

\noindent We need the next lemma linking sources of simple modules to the structure of the Dade group, the proof of which consists of the arguments given in \cite[paragraph after Proposition~2.1]{CELK11}. 


\begin{lem}\label{lem:SourcesSimples_p_soluble_groups}\label{lem:sourcesSimpleDihedral}\label{lem:SourcesSimplesSolubleGroups}
Let $G$ be a $p$-soluble group. Let $T$ be a simple $kG$-module, let $Q$ be a vertex of~$T$ and let $U$ be a $kQ$-source of $T$.  Then, the following assertions hold:
    \begin{enumerate}[label={\rm(\alph*)}]
        \item $U$ is a capped endo-permutation $kQ$-module and its class in the Dade group is a torsion element, i.e. $[U]\in\bD_k^{tors}(Q)$; and 
        \item if $p=2$ and  $Q\cong D_{2^n}$ with $n\geq 3$, then $U\cong k$. 
    \end{enumerate}
\end{lem}    
%\todo[color=green]{
%For $p=2$:\\$2$-soluble= soluble. We apply this result to $2$-nilpotent groups.}
    
\begin{proof}    
    (a) First, $U$ is an endo-permutation module by \cite{PuigNilpotentBlocks88}. Furthermore, by a result of Feit, the simple module $T$ is algebraic; see \cite{Feit80}. Thus, \cite[Proposition~4.4]{BoltjeKuls06} proves that the class of $U$ in the Dade group $\bD_k(Q)$ must be a torsion element.\\
    (b) As $Q$ is a dihedral $2$-group, we have  $\bD_k^{tors}(Q)=\{[k]\}$ (see e.g. \cite[Proposition~6.4(d)]{ThevenazSurvey}), proving that $U$ is the trivial $kQ$-module. 
\end{proof}




\vspace{2mm}
\subsection{Clifford-theoretic results}

We will use the standard reduction techniques due to 
Fong and Reynolds. We review them here in terms of splendid and endo-permutation source Morita equivalences, which is less standard in the literature. These techniques work irrespective of the characteristic, so in this subsection $p$ is an arbitrary prime.


\begin{thm}[Fong{--}Reynolds, 
see {\cite[1.5. Theorem]{KoshitaniKunugiWaki04}}]\label{thm:FongReynolds}%
Let $G$ be a finite group, let $N\unlhd G$, let $\bb\in\Bl(kN)$
and  let $T:= T_G(\bb)$ be the stabilizer of $\bb$ in $G$. Then, block induction induces a bijection
$$\Bl(kT\mid\bb)\longrightarrow \Bl(kG\mid\bb), \bA\mapsto \bA^G\,.$$
Moreover, if $\bB\in\Bl(kG\mid\bb)$ and $\widetilde{\bB}\in\Bl(kT\mid\bb)$ is the preimage of $\bB$ under the latter bijection, then the pair of bimodules $(e^{}_{\bB}\cdot kG\cdot e^{}_{\widetilde{\bB}}, e^{}_{\widetilde{\bB}}\cdot kG\cdot e^{}_{\bB})$ realizes a splendid Morita equivalence between $\bB$ and $\widetilde{\bB}$. 
\end{thm}

\begin{thm}[Fong’s Second Reduction Theorem, see {\cite[Theorem~6.8.13]{LinckelmannVol2}} and  {\cite[ Proposition~3.3]{SalminenSOURCES08}}]\label{thm:Fong2ndReduction}%
   Let $G$ be a finite group and let $N\unlhd G$ be a normal subgroup. Let $\bB\in\Bl(kG)$ be a block  covering a $G$-stable block $\bc\in\Bl(kN)$ and set $S:=kNe_{\bc}$. 
   If $D$ is a defect group of $\bB$ and $D\cap N=1$, then $\bc$ has defect zero and there is an isomorphism of interior $D$-algebras
   \[
   kGe_{\bc}\cong kNe_{\bc}\otimes_k k^{}_{\alpha}[G/N]
   \]
   where $\alpha\in {\mathrm{H}}^2(G/N,k^{\times})$.  Moreover, the following assertions hold.
   \begin{enumerate}[label={\rm(\alph*)}]
       \item  
       There exists a central extension 
       {$1\longrightarrow Z\longrightarrow H\longrightarrow G/N\longrightarrow 1$}, where $Z$ is the cyclic group of $|G/N|_{p^\prime}$-th roots of unity in $k$, and there exists an idempotent $e\in kZ$ such that $k^{}_{\alpha}[G/N]\cong kHe$.
       \item The block $\bB$ is Morita equivalent to a block $\widehat{\bB}$ of $kH$ having a defect group $\widehat D$ isomorphic to~$D$. After identifying these defect groups, the Morita equivalence is induced by a bimodule with vertex $\Delta D$ and an endo-permutation source.
       \item  If $\alpha$ is trivial, then $H$ can be replaced by $G/N$ itself. 
       If $G/N$ is perfect, then $H$ can be replaced by $[H,H]$ and $Z$ by $Z([H,H])$. 
   \end{enumerate}   
\end{thm}

\noindent Detailed proofs of the main claim, of assertion (a) and of the second claim of assertion  (c)  can be found in \cite[\S3 and \S4]{SalminenSOURCES08}.
    The first claim of assertion (c) follows from \cite[Proposition 1.2.6]{LinckelmannVol1}.
        %See also \cite[Proposition~1.2.18]{LinckelmannVol1} for the replacement of the twisted group algebra by $kHe$. 
   Assertion~{\rm(b)}, and in particular the more precise information on the vertex and source of the bimodule inducing the Morita equivalence, will be crucial in our subsequent arguments. We record this information in the
following remark. Its proof, together with a proof of Assertion~{\rm(b)}, is given in
Appendix~\ref{app:A}.

\begin{rem}\label{rem:ProofFong2ndRed}%
    %\begin{enumerate}[label=(\arabic*)]
        %\item
       % \item  
       Since $\bc$ has defect zero, the block algebra $S:=kNe_{\bc}$ is a matrix algebra. Thus, $S\cong\End_k(V)$ where $V$ is, up to isomorphism, the unique simple $\bc$-module. The module $V$ has a unique extension, up to isomorphism, to a simple $k[ND]$-module, which we also denote by $V$. This extension has vertex $D$, and $\Res_D^{ND}(V)$ is a capped endo-permutation $kD$-module. Moreover, viewed as a $k[\Delta D]$-module via the identification $D\cong\Delta D$, the cap of $\Res_D^{ND}(V)$ is a source of the bimodule inducing the Morita equivalence constructed in the proof of Assertion~{\rm(b)}.
   % \end{enumerate}
\end{rem}




\vspace{2mm}
\subsection{Splendid stable equivalences of Morita type and trivial source modules}
The following lemma generalizes parts of~\cite[Lemma A.3]{KMN11}. 

\begin{lem}\label{lem:stablesplendid_and_pperm_modules}
     Let $G$ and $H$ be two subgroups of a finite group $L$. Assume  $G$ and $H$ have a common Sylow $p$-subgroup $P$ and the same fusion system on $P$, i.e. $\mathcal{F}_P(G)=\mathcal{F}_P(H)$.
     Let $\bA\in\Bl(kG)$ and let $\bB\in\Bl(kH)$ be blocks with defect group $P$. Assume
     $M$ (resp. $M^{\prime}$) is an indecomposable  finitely generated $(\bA,\bB)$-bimodule (resp. a $(\bB,\bA)$-bimodule),  which,  
     viewed as a $k[G\times H]$-module (resp. as a $k[H\times G]$-module), is  a $p$-permutation module with vertex $\Delta P$, and assume that the pair $(M, M^{\prime})$ induces a stable equivalence of Morita type between $\bA\mathrm{-mod}$ and $\bB\mathrm{-mod}$.
     \par
     Assume, furthermore, that $X$ is a non-projective indecomposable $\bB$-module. Then, the following assertions hold.
     \begin{enumerate}[label={\rm(\alph*)}]
         \item There exists a non-projective indecomposable $\bA$-module $Y$, unique up to isomorphism, such that $M\otimes_{\bB}X = Y\oplus (\mathrm{proj})$ as $\bA$-modules. 
         \item If $X$ is a trivial source module, then so is $Y$. Moreover, $X$ and $Y$ have a vertex in common.
     \end{enumerate}
\end{lem}

\begin{proof}
    Assertion~(a) is well-known: it  follows from \cite[Proposition 4.13.2]{LinckelmannVol1}. 
    We prove (b). As $X$ is non-projective, we may choose a vertex $Q$ of $X$ such that $1\neq Q\leq P$ and we have 
    \[    
    X \big| \Ind_Q^H(k)=kH\otimes_{kQ}k\,.
    \]
    Now, by the assumption, we have $M\mid \Ind_{\Delta P}^{G\times H}(k)$, which is $kG\otimes_{kP}kH$ as  $(kG,kH)$-bimodule.
    Therefore, 
    \[
    Y\,\big| \left(M\otimes_{kH}X\right) \big|\,  kG\otimes_{kP}kH \otimes_{kH}kH\otimes_{kQ}k \,\big|\, \Ind_P^G\left(   
                \Res^H_P(\Ind_{Q}^{H}(k))
            \right)
   \]
   and the Mackey Formula yields 
   \[
    Y \,\big|   \bigoplus_{g\in [Q\backslash G/P]}\Ind_{\lconj{g}{Q}\cap P}^{G}(k)\,.
   \]
   Hence, $Y$ is a trivial source module, and it has a vertex $R$ such that  $R\leq\!\lconj{g}{Q}\cap P$ for some $g\in G$.  It follows that $\lconj{g^{-1}}{R}\leq Q\cap \lconj{g^{-1}}{P}\leq Q\leq P$. Since $\mathcal{F}_P(G)=\mathcal{F}_P(H)$, there exists $h\in H$ such that   $\lconj{g^{-1}}{R}=\lconj{h}{R}$, and so $\lconj{h}{R}$ is also a vertex of $Y$ and $R\leq_H Q$. Since the pair $(M, M^{\prime})$ induces a stable equivalence of Morita type, the same argument with $M^{\prime}$ yields $Q\leq_G R$. It follows that $R=_{H}Q$ and $Q=_GR$. Assertion (b) follows. 
\end{proof}
